编译 LaTeX 经常会遇到报错，以下是常见错误类型及排查手段的速查表：

\begin{table}[htbp]
\centering
\caption{LaTeX 排错速查表}
\begin{tabular}{lp{10cm}}
\hline
\textbf{错误提示} & \textbf{原因与解决办法} \\
\hline
\texttt{Undefined control sequence} & 拼写错误或未引入对应的宏包。请检查命令拼写，或在导言区添加所需的 \texttt{\textbackslash usepackage\{\}}。 \\
\texttt{File not found} & 路径或文件名错误。检查包含文件或图片的路径，确保使用正斜杠 \texttt{/} 而非反斜杠。 \\
\texttt{Missing \$ inserted} & 数学公式环境错误。检查公式中是否混入了普通文本，或者行内公式遗漏了 \texttt{\$} 符号。 \\
\texttt{Package xxx Warning} & 宏包警告，不影响最终生成 PDF，通常可以忽略。若有强迫症可根据提示调整参数。 \\
\hline
\end{tabular}
\end{table}

\subsection{通用自救指南}
\begin{enumerate}
    \item 仔细阅读终端输出的 \texttt{Error} 报错行，通常会指明报错的具体行号。
    \item 若无法定位问题，尝试将近期新增的代码注释掉，重新编译，以二分法定位错误源。
    \item 若更改了核心结构或 \texttt{.bib} 参考文献，务必清理所有的辅助文件（如 \texttt{.aux}, \texttt{.bbl}, \texttt{.log} 等）后再次完整编译。
\end{enumerate}
